% !TeX program = xelatex
% Μεταγλώττιση: xelatex lec1_exercises.tex
% Στο Overleaf: επιλέξτε XeLaTeX ως compiler.
\documentclass[11pt,a4paper]{article}

\usepackage[margin=1.9cm]{geometry}
\usepackage{fontspec}
\usepackage{polyglossia}
\setdefaultlanguage{greek}
\setotherlanguage{english}
\setmainfont{DejaVu Serif}[ItalicFont={DejaVu Serif},ItalicFeatures={FakeSlant=0.2},BoldItalicFont={DejaVu Serif Bold},BoldItalicFeatures={FakeSlant=0.2}]
\setmonofont{DejaVu Sans Mono}[Scale=0.88]
\newfontfamily\greekfont{DejaVu Serif}[ItalicFont={DejaVu Serif},ItalicFeatures={FakeSlant=0.2},BoldItalicFont={DejaVu Serif Bold},BoldItalicFeatures={FakeSlant=0.2}]
\newfontfamily\greekfonttt{DejaVu Sans Mono}[Scale=0.88]
\usepackage{amsmath,amssymb,amsthm}
\usepackage{enumitem}
\usepackage{fancyvrb}
\usepackage{fancyhdr}
\usepackage[unicode,hidelinks]{hyperref}
\hypersetup{pdftitle={Το πρόβλημα του τερματισμού και τρεις μαθηματικές εικασίες}}
\newcommand{\horrule}[1]{\rule{\linewidth}{#1}}
\theoremstyle{definition}
\newtheorem{question}{Άσκηση}
\newcommand{\HALT}{\operatorname{HALT}}
\setlist[enumerate]{label=(\alph*),leftmargin=*,itemsep=3pt,topsep=5pt}
\setlength{\parindent}{0pt}
\setlength{\parskip}{5pt}
\setlength{\emergencystretch}{3em}
\pagestyle{fancy}
\fancyhf{}
\fancyfoot[L]{\small ΗΥ280 (2026) - Τσουρακάκης}
\fancyfoot[R]{\small\thepage}
\renewcommand{\headrulewidth}{0pt}
\DefineVerbatimEnvironment{pseudocode}{Verbatim}{fontsize=\small,xleftmargin=1em}

\begin{document}
\begin{center}
\horrule{0.4pt}\\[0.15cm]
{\large\textbf{ΗΥ 280 \textemdash{} Θεωρία Υπολογισμού}}\\[0.1cm]
{Ασκήσεις 2ης  διάλεξης}\\[0.1cm]
{Χαράλαμπος Ε. Τσουρακάκης}\\
\horrule{0.4pt}
\end{center}

\begin{center}
\fbox{%
\begin{minipage}{0.93\linewidth}
\textbf{Οδηγίες για τις ασκήσεις}

\begin{itemize}[leftmargin=*,itemsep=6pt]
    \item \textbf{Ασκήσεις 1-4:}
    Αποτελούν μέρος του παραδοτέου των ασκήσεων
    του μαθήματος.  

    \item \textbf{Άσκηση 4 ($\bigstar \bigstar \bigstar$):}
    Είναι απαιτητική. Ασχοληθείτε μαζί της
    μόνο αφού έχετε ολοκληρώσει και κατανοήσει
    τις τρεις πρώτες ασκήσεις. 


    \item \textbf{Άσκηση 5 — προαιρετική:}
 Δεν αποτελεί μέρος του παραδοτέου, αλλά μέρος της περιέργειας που χαρακτηρίζει έναν επιστήμονα υπολογιστών: μπορεί ένα πρόγραμμα να εκτυπώσει τον ίδιο του τον κώδικα; Η άσκηση είναι αρκετά απαιτητική, ιδιαίτερα αν συναντάτε για πρώτη φορά την έννοια της αυτοαναφοράς. Χρειάζεται ένα ευρηματικό τέχνασμα, οπότε μην απογοητευτείτε αν δεν το ανακαλύψετε αμέσως! Αν δεν νιώθετε ακόμη άνετα με κάποια γλώσσα προγραμματισμού, προτείνω να την αφήσετε.

     
    \item {\bf Αξιοποιήστε τις ώρες γραφείου} για τις απορίες σας, αφού κατανοήσετε τις ασκήσεις και τις προσπαθήσετε.  
\end{itemize}
\end{minipage}%
}
\end{center}
\medskip


\begin{question}[$\bigstar$ Διαγωνιοποίηση και προτιμήσεις παγωτού]
Τρεις μαθητές δηλώνουν ποιες γεύσεις παγωτού τούς αρέσουν:
\begin{center}
\renewcommand{\arraystretch}{1.3}
\begin{tabular}{lccc}
\hline
 & Σοκολάτα & Βανίλια & Φράουλα \\
\hline
James & Ναι & Όχι & Ναι \\
Joe   & Όχι & Όχι & Ναι \\
John  & Ναι & Ναι & Ναι \\
\hline
\end{tabular}
\end{center}
Ένα \emph{προφίλ προτιμήσεων} είναι μια τριάδα απαντήσεων
«Ναι» ή «Όχι», μία για κάθε γεύση με την παραπάνω σειρά.

\begin{enumerate}
\item Κατασκευάστε ένα προφίλ που διαφέρει από εκείνο
του James στη σοκολάτα, του Joe στη βανίλια
και του John στη φράουλα.

\item Εξηγήστε γιατί το προφίλ σας είναι διαφορετικό
από καθένα από τα τρία δοσμένα προφίλ.
Χρειάζεται να διαφέρει από κάθε μαθητή και στις τρεις γεύσεις;

\item Κωδικοποιήστε το «Ναι» με $1$ και το «Όχι» με $0$.
Αν $a_{ij}$ είναι η απάντηση του $i$-οστού μαθητή
για την $j$-οστή γεύση, δώστε έναν τύπο
για την $i$-οστή απάντηση $b_i$ του νέου προφίλ.

\item Γενικεύστε: δίνονται $n$ μαθητές και $n$ γεύσεις,
όπου $n\ge1$.
Αποδείξτε ότι μπορείτε πάντοτε να κατασκευάσετε
ένα προφίλ διαφορετικό από τα προφίλ όλων των μαθητών,
ακόμη και αν κάποιοι μαθητές έχουν ίδιες προτιμήσεις.

\end{enumerate}
\end{question}



\begin{question}[ $\bigstar\bigstar $ Το διαγώνιο επιχείρημα του Cantor]
Στόχος της άσκησης είναι να αποδείξετε ότι το διάστημα
$(0,1)$ είναι μη αριθμήσιμο.

Υποθέτουμε προς άτοπο ότι υπάρχει μια λίστα
$x_1,x_2,x_3,\ldots$ που περιέχει όλους τους αριθμούς
του $(0,1)$. Γράφουμε:
\[
\begin{array}{rcl}
x_1 &=& 0{,}a_{11}a_{12}a_{13}a_{14}\cdots \\
x_2 &=& 0{,}a_{21}a_{22}a_{23}a_{24}\cdots \\
x_3 &=& 0{,}a_{31}a_{32}a_{33}a_{34}\cdots \\
\vdots && \vdots
\end{array}
\]
όπου $a_{ij}\in\{0,1,\ldots,9\}$.

Μπορείτε να χρησιμοποιήσετε το εξής γνωστό γεγονός:
κάθε πραγματικός αριθμός στο $(0,1)$ έχει μοναδική
δεκαδική ανάπτυξη που δεν τελειώνει σε άπειρα $9$.
Για κάθε $x_i$ επιλέγουμε αυτή την ανάπτυξη.

\begin{enumerate}
\item Τι συμβολίζει το $a_{ij}$;
Ποια ψηφία βρίσκονται στη διαγώνιο του παραπάνω πίνακα;

\item Γιατί αρκεί να κατασκευάσουμε έναν αριθμό
$y\in(0,1)$ τέτοιο ώστε $y\ne x_i$ για κάθε $i\ge1$;

\item Θέλουμε να κατασκευάσουμε τον αριθμό
\[
y=0{,}b_1b_2b_3\cdots
\]
χρησιμοποιώντας μόνο τα ψηφία $1$ και $2$.
Συμπληρώστε τον κανόνα ώστε $b_i\ne a_{ii}$:
\[
b_i=
\begin{cases}
\underline{\hspace{8mm}}, & \text{αν }a_{ii}=1,\\[2mm]
\underline{\hspace{8mm}}, & \text{αν }a_{ii}\ne1.
\end{cases}
\]

\item Αν τα πρώτα πέντε ψηφία της διαγωνίου είναι
$1,4,2,1,0$, ποια είναι τα πρώτα πέντε ψηφία του $y$;

\item Αποδείξτε ότι $y\in(0,1)$.

\item Γνωρίζουμε ότι
\[
0{,}5000\ldots=0{,}4999\ldots
\]
Γιατί η διαφορά σε ένα δεκαδικό ψηφίο δεν αρκεί
γενικά για να συμπεράνουμε ότι δύο αριθμοί διαφέρουν;
Εξηγήστε γιατί το πρόβλημα αυτό δεν εμφανίζεται
στη συγκεκριμένη κατασκευή.

\item Αποδείξτε ότι $y\ne x_i$ για κάθε $i\ge1$
και ολοκληρώστε την απόδειξη.
Χρειάζεται ο $y$ να διαφέρει από τον $x_i$
σε όλα τα ψηφία;

\item Κάποιος προτείνει:
«Ας προσθέσουμε τον $y$ στην αρχή της λίστας.
Τώρα η λίστα είναι πλήρης!»
Εξηγήστε γιατί η πρόταση αυτή δεν επιλύει το πρόβλημα.
\end{enumerate}
\end{question}




\begin{question}[$\bigstar \bigstar$ Η λογική μορφή της διαγωνιοποίησης]
Έστω μη κενό σύνολο $S$.
Μια συνάρτηση $f:\mathbb N_{>0}\to S$ περιγράφει τη λίστα
\[
f(1),f(2),f(3),\ldots
\]
Επιτρέπονται επαναλήψεις στοιχείων.

\begin{enumerate}
\item Για δεδομένη συνάρτηση $f$, γράψτε με ποσοδείκτες:
«Η λίστα που περιγράφει η $f$ περιέχει κάθε στοιχείο του $S$».

\item Χρησιμοποιώντας το προηγούμενο ερώτημα,
διατυπώστε συμβολικά:
«Υπάρχει πλήρης λίστα των στοιχείων του $S$».

\item Αρνηθείτε την πρόταση του προηγούμενου ερωτήματος.
Μεταφέρετε την άρνηση προς τα μέσα, αλλάζοντας διαδοχικά
τους ποσοδείκτες, μέχρι να εμφανιστεί μόνο μια ανισότητα
στο εσωτερικό.
Σε κάθε βήμα αναφέρετε τον κανόνα που χρησιμοποιείτε.

\item Με βάση τον τελικό τύπο, συμπληρώστε:
\begin{itemize}
    \item Μας δίνεται μια αυθαίρετη \underline{\hspace{25mm}}.
    \item Πρέπει να βρούμε ένα \underline{\hspace{25mm}}.
    \item Πρέπει να δείξουμε ότι, για κάθε $i\ge1$,
          \underline{\hspace{35mm}}.
\end{itemize}
Το στοιχείο που βρίσκουμε επιτρέπεται να εξαρτάται
από τη λίστα $f$; Αιτιολογήστε.

\item Ένας φοιτητής γράφει:
\[
\exists s\in S\;
\forall f:\mathbb N_{>0}\to S\;
\forall i\ge1:\ s\ne f(i).
\]
Τι ισχυρίζεται αυτός ο τύπος;
Γιατί είναι ψευδής για κάθε μη κενό $S$;

Υπόδειξη: για ένα υποψήφιο $s$, εξετάστε
τη σταθερή λίστα $f(i)=s$.
\end{enumerate}
\end{question}





\begin{question}[$\bigstar\bigstar\bigstar$ Διαγωνιοποίηση και σταθερά σημεία]
Έστω $T$ ένα μη κενό σύνολο και $f:T\to T$ μια συνάρτηση
\emph{χωρίς σταθερά σημεία}, δηλαδή
\[
f(t)\ne t \qquad \text{για κάθε }t\in T.
\]
Αποδείξτε ότι, για κάθε μη κενό σύνολο $S$, δεν υπάρχει
επί συνάρτηση
\[
g:S\longrightarrow T^S,
\]
όπου
\[
T^S=\{h\mid h:S\to T\}
\]
είναι το σύνολο όλων των συναρτήσεων από το $S$ στο $T$.
\end{question}



\begin{question}[$\bigstar\bigstar\bigstar\bigstar$ Ένα πρόγραμμα που εκτυπώνει τον εαυτό του]
 

Στη γλώσσα προγραμματισμού της επιλογής σας,
γράψτε ένα \emph{quine}: ένα πρόγραμμα του οποίου
η έξοδος είναι ακριβώς ο ίδιος ο πηγαίος κώδικάς του.

Το πρόγραμμα δεν πρέπει να δέχεται είσοδο ούτε
να διαβάζει το αρχείο του πηγαίου κώδικά του.
Προσέξτε τα εισαγωγικά, τα κενά και τις αλλαγές γραμμής.

Εξηγήστε πώς το πρόγραμμά σας επιτυγχάνει
την αυτοαναπαραγωγή του χωρίς να διαβάζει τον εαυτό του
από κάποιο αρχείο.
\end{question}



\end{document}
